\documentclass[12pt,a4paper]{article}

\usepackage[czech]{babel}
\usepackage[T1]{fontenc}
\usepackage[utf8]{inputenc}

\usepackage{geometry}
\usepackage[most]{tcolorbox}

\geometry{margin=2.5cm}

\title{\textbf{Rozbor díla „Kulička“}}
\author{Pavlo}
\date{}

\begin{document}
\maketitle

%------------------------------------------------
\section*{Autor}

\begin{tcolorbox}[breakable]

\textbf{Autor:}
\begin{itemize}
\item Guy de Maupassant
\end{itemize}

\textbf{Název díla:}
\begin{itemize}
\item Kulička
\end{itemize}

\textbf{Literární směr:}
\begin{itemize}
\item Realismus a naturalismus
\end{itemize}

\textbf{Století tvorby:}
\begin{itemize}
\item 2. polovina 19. století
\end{itemize}

\textbf{Další autoři stejného směru + dílo:}
\begin{itemize}
\item Émile Zola – Zabiják
\item Gustave Flaubert – Paní Bovaryová
\end{itemize}

\textbf{Další autorova díla:}
\begin{itemize}
\item Slečna Fifi
\item Miláček
\item Mont-Oriol
\end{itemize}

\textbf{Specifický styl autora:}
\begin{itemize}
\item Autor často zachycuje realistické situace ze života společnosti, zaměřuje se na mezilidské vztahy a pokrytectví. Často využívá prostředí prusko-francouzské války.
\end{itemize}

\end{tcolorbox}

%------------------------------------------------
\section*{Charakteristika uměleckého textu}

\begin{tcolorbox}[breakable]

\textbf{Literární druh, žánr, forma:}
\begin{itemize}
\item Epika, povídka, próza
\end{itemize}

\textbf{Typ vypravěče:}
\begin{itemize}
\item Er-forma, vševědoucí vypravěč (heterodiegetický)
\end{itemize}

\textbf{Dominantní slohový postup:}
\begin{itemize}
\item Vypravovací s prvky popisu
\end{itemize}

\textbf{Vysvětlení názvu díla:}
\begin{itemize}
\item Název je odvozen od přezdívky hlavní postavy, která souvisí s jejím kulatým vzhledem.
\end{itemize}

\textbf{Aktuálnost díla:}
\begin{itemize}
\item Dílo je aktuální i dnes, protože ukazuje pokrytectví společnosti a rozdíly mezi společenskými vrstvami.
\end{itemize}

\textbf{Místo a čas děje:}
\begin{itemize}
\item Normandie, rok 1870, období prusko-francouzské války
\end{itemize}

\textbf{Stručné nastínění děje:} \\
Děj se odehrává během prusko-francouzské války. Skupina cestujících z různých společenských vrstev opouští okupovaný Rouen a vydává se dostavníkem do Dieppe. Mezi nimi je i prostitutka Alžběta Roussetová, přezdívaná Kulička.\\
\\
Zpočátku ji ostatní opovrhují, ale když se s nimi rozdělí o jídlo, začnou ji přijímat. V hostinci je však jejich cesta zastavena pruským důstojníkem, který požaduje, aby se mu Kulička oddala. Ta to z vlasteneckých důvodů odmítá.\\
\\
Ostatní cestující ji postupně přemlouvají, aby se obětovala pro jejich dobro. Nakonec Kulička ustoupí. Po odjezdu ji však všichni opět odmítají a chovají se k ní s pohrdáním.\\

\textbf{Smysl díla:}
\begin{itemize}
\item Kritika pokrytectví společnosti a morálních rozdílů mezi lidmi různých vrstev.
\end{itemize}

\textbf{Pravděpodobný adresát:}
\begin{itemize}
\item Čtenáři 19. století, především francouzská společnost
\end{itemize}

\textbf{Kompozice:}
\begin{itemize}
\item Chronologická
\end{itemize}

\textbf{Zařazení do kontextu autorovy tvorby:}
\begin{itemize}
\item Jedno z nejznámějších a nejvýznamnějších děl autora
\end{itemize}

\textbf{Film nebo dramatizace:}
\begin{itemize}
\item Existují filmová zpracování, např. film „Kulička“ z roku 1934
\end{itemize}

\end{tcolorbox}

%------------------------------------------------
\section*{Úryvek k rozboru}

Byla to žena kulatá, malá, s prsty jako párky, s lesklou pletí a s černýma očima, které se zdály být stále plné smíchu. Její tvář byla svěží jako jablko a její postava připomínala malou kouli, odtud také pocházela její přezdívka Kulička. Její šaty byly vždy upravené, i když prosté, a bylo na ní vidět, že se snaží působit důstojně. Když mluvila, její hlas byl měkký a příjemný, ale občas se v něm objevil nádech rozhodnosti. Ostatní cestující se na ni dívali s určitým opovržením, přestože si nemohli nevšimnout její dobrosrdečnosti. V jejím chování bylo cosi upřímného, co kontrastovalo s chladnou zdvořilostí ostatních. Když se usmála, zdálo se, že kolem ní vše ožívá. Byla to žena, která se lišila od ostatních nejen vzhledem, ale i srdcem.

\begin{tcolorbox}[breakable]

\textbf{Atmosféra úryvku:}
\begin{itemize}
\item Atmosféra je klidná a popisná, zároveň však jemně napjatá díky kontrastu mezi dobrosrdečností hlavní postavy a odmítavým postojem ostatních.
\end{itemize}

\textbf{Postavy v úryvku:}
\begin{itemize}
\item Kulička (Alžběta Roussetová)
\item ostatní cestující (měšťané, šlechta)
\end{itemize}

\textbf{Charakteristika postav:}

\textbf{Hlavní postava}
\begin{itemize}
\item Kulička – dobrosrdečná, upřímná, laskavá, společensky odsuzovaná, ale morálně vyspělejší než ostatní
\end{itemize}

\textbf{Vedlejší postavy}
\begin{itemize}
\item měšťané – pokrytečtí, sobečtí, povýšení
\item šlechta – chladná, odtažitá, zdůrazňuje společenské postavení
\end{itemize}

\textbf{Další postavy díla:}
\begin{itemize}
\item manželé Loiseauovi
\item manželé Carré-Lamadonovi
\item hrabě a hraběnka
\item dvě jeptišky
\item Cornudet
\end{itemize}

\textbf{Vztahy mezi postavami:}
\begin{itemize}
\item ostatní postavy Kuličkou pohrdají, ale zároveň ji využívají
\item Kulička je k nim zpočátku vstřícná a obětavá
\end{itemize}

\textbf{Zařazení úryvku do děje:}
\begin{itemize}
\item úryvek se nachází na začátku díla a představuje hlavní postavu
\end{itemize}

\textbf{Jazykové prostředky:}
\begin{itemize}
\item přirovnání
\item kontrast
\item metafora
\item ironie
\end{itemize}

\end{tcolorbox}

%------------------------------------------------
\section*{Názor na dílo}

Dílo Guy de Maupassanta poutavě zachycuje vztahy mezi společenskými vrstvami a poukazuje na jejich pokrytectví. I přes dobu vzniku se čte poměrně snadno a obsahuje výstižné popisy postav. Zároveň ukazuje, jak válka ovlivňuje chování lidí.\\
\\
Dílo bych doporučil.

\end{document}
